{\color{orange}

\section*{Interesting Questions for a DAOS Benchmarking Paper}

\subsection*{Performance Characterization}

\begin{itemize}
    \item How does DAOS's userspace I/O stack compare to kernel-bypass alternatives (SPDK-direct, NVMe-oF) at various concurrency levels?
    \item Where exactly does the bottleneck shift as you scale --- PMem bandwidth, NVMe throughput, network fabric, or DAOS engine threads?
    \item What is the overhead of DAOS's epoch/transaction model vs.\ raw device performance?
\end{itemize}

\subsection*{Workload Sensitivity}

\begin{itemize}
    \item How does the 4K--1M I/O size sweep expose DAOS's tiering decisions between PMem and NVMe?
    \item Does the read/write asymmetry change character across workload types (random, sequential, mixed)?
    \item How does object-level access (DAOS native API) compare to POSIX via DFS at the same block size?
\end{itemize}

\subsection*{Scalability and Concurrency}

\begin{itemize}
    \item Does performance degrade gracefully with many clients, or are there cliff edges?
    \item How does server-side thread count interact with client count --- is it a known bottleneck?
\end{itemize}

\subsection*{Practical HPC Relevance}

\begin{itemize}
    \item Does the benchmark reflect realistic application I/O patterns (checkpoint, analysis, streaming)?
    \item What is the gap between peak micro-benchmark numbers and what a real MPI application sees?
\end{itemize}

\subsection*{Core Question}

The question worth pushing hardest: \textbf{where does DAOS's architectural novelty
(PMem-first, userspace, asynchronous) actually show up in the numbers, vs.\ where it is
indistinguishable from a well-tuned Lustre or GPFS setup?}
That is the paper's core claim to make or break.

}
